\chapter{Теория}

\section{Логические методы классификации}
Логические методы классификации — это методы машинного обучения, которые принимают решение на основе логических правил вида "если-то". \\

Алгоритм строит дерево решений, где каждый узел — проверка условия по признаку,
ветви — результат проверки (True / False),
листья — итоговый класс.


\subsection{Отличие от метрических методов}
Метрические методы классифицируют объект на основе расстояния между объектами в пространстве признаков. Объект относится к тому классу, к которому принадлежат ближайшие объекты обучающей выборки.

\section{max\_depth и max\_features}

max\_depth — это параметр, который ограничивает максимальную глубину дерева решений. \\

max\_features — это параметр, который определяет максимальное количество признаков, которые рассматриваются при поиске лучшего разбиения узла.


\endinput